\section{Experimental Setup}
\label{sec:setup}

\textbf{Datasets.} \Cref{tab:datasets} summarises the three datasets. All are replayed in timestamp order.
From each we take 400{,}000 flows evenly spaced over time (deterministic); the first 30\% form the training
period (its last 20\% is the reference set) and the remaining 280{,}000 flows form 70 windows of 4{,}000 flows.
CIC-IDS2017~\cite{sharafaldin2018} uses the \texttt{TrafficLabelling} CSVs, whose timestamps lack AM/PM; we
restore afternoon times. The fitted part of its training period (Monday and Tuesday morning) contains only
benign traffic and FTP-Patator, so twelve attack classes are new. For CSE-CIC-IDS2018 we use the ``Processed Traffic Data for ML''
CSVs, thinned to one row in eight per file before sampling; seven attack classes are known at training time and
seven appear later. In UNSW-NB15~\cite{moustafa2015unsw} every class already occurs in the first capture day, so,
as is common in open-set NIDS work, we remove DoS, Reconnaissance, Shellcode and Worms from the training
period; they appear in the stream as unknown attacks.

\begin{table}[t]
\centering
\caption{Datasets (after replay sampling: 400k flows, 70 windows).}
\label{tab:datasets}
\resizebox{\columnwidth}{!}{%
\begin{tabular}{@{}lccc@{}}
\toprule
 & CIC-IDS2017 & CSE-CIC-IDS2018 & UNSW-NB15 \\
\midrule
Source flows & 2.83\,M & 16.2\,M & 2.54\,M \\
Days covered & 5 & 10 & 2 \\
Classes (incl.\ benign) & 15 & 15 & 10 \\
Attack classes in training & 1 & 7 & 5 \\
New in stream & 12 & 7 & 4 (held out) \\
\bottomrule
\end{tabular}}
\end{table}

\textbf{Labels.} Budget $B = 10\%$ (main) and $1\%$ of each window, delay $\delta = 1$ window. Every method that
uses labels gets exactly the same number; INSOMNIA uses none by design.

\textbf{Model.} All methods use the same classifier: scikit-learn histogram gradient boosting (150 iterations,
L2 regularisation 1.0); classes with fewer than five labelled flows are left out of training and stay \unk{}.
All methods also use the same open-world detector (maximum softmax probability and class-conditional $k$-NN),
except in the detector ablation, so differences come from the adaptation strategy.

\textbf{Baselines.} \emph{Static} (never updated); \emph{periodic retrain} (from scratch every five windows on
the training data and the last five labelled windows); \emph{continual replay} (uniform reservoir memory, retrain
whenever labels arrive); \emph{guarded selective} (drift- or novelty-triggered retrain with a hold-out guard and
rollback). Re-implemented under our protocol: \emph{uncertainty sampling} (least-confidence query, INSOMNIA's
``US + oracle'' upper bound~\cite{andresini2021insomnia}); \emph{INSOMNIA}~\cite{andresini2021insomnia}
(least-confident 50\% of each window labelled by a nearest-centroid estimator, no analyst labels);
\emph{NetGuard PS}~\cite{gupta2025netguard} (density-ratio query with two diagonal Gaussian mixtures; the
generative augmentation is not reproduced); and \emph{CBRS replay}~\cite{chrysakis2020online,amalapuram2023augmented}
(class-balancing reservoir memory). All replay memories hold 20{,}000 flows.

\textbf{Development and held-out data.} All design decisions, including the ablation, were made on
CIC-IDS2017. CSE-CIC-IDS2018 and UNSW-NB15 were used for no design decision (we only checked their class
composition with one run of the static model, which led to the UNSW-NB15 hold-out above). Both are
evaluated only with the configuration frozen on CIC-IDS2017.

\textbf{\nova{} hyper-parameters} were fixed from design arguments and two development runs, then frozen:
uniform share $\rho_u = 0.2$, $K = 24$ clusters, $k_p = 5$, $r = 5$, benign weight $\beta = 4$, 25\% hold-out,
calibration store of 200 flows per class (after 50 in memory), minimum support $m = 5$, $\epsilon = 0.05$,
rollback margin $0.05$, recency weight $\omega = 2$, $k = 10$ with 3{,}000 flows per class bank.

\textbf{Runs.} Three seeds per configuration (the seed changes model initialisation, label sampling and
clustering, not the data); we report mean $\pm$ standard deviation. To compare two methods we take the
per-window difference of their seed-averaged OW-F1 (or FPR) and report its mean with a 95\% moving-block bootstrap
interval~\cite{kunsch1989jackknife} (blocks of five windows, 10{,}000 resamples), which respects the autocorrelation of consecutive windows;
we call a difference significant only if the interval excludes zero. Neural variants use PyTorch 2.14.1 with CUDA
13.0 on an RTX 3060; everything else runs on a 16-core CPU.
